\documentclass[runningheads]{llncs}

\usepackage[T1]{fontenc}
\usepackage{graphicx}
\usepackage{amsmath}
\usepackage{booktabs}
\usepackage{float}
\usepackage[hidelinks]{hyperref}
\usepackage{array}
\usepackage{cite}
\usepackage{bm}

\setlength{\textfloatsep}{8pt plus 2pt minus 2pt}
\setlength{\floatsep}{8pt plus 2pt minus 2pt}
\setlength{\intextsep}{8pt plus 2pt minus 2pt}
\setlength{\abovecaptionskip}{4pt}
\setlength{\belowcaptionskip}{2pt}


\begin{document}

\newcommand{\best}[1]{\bm{#1}}

\title{Solar Flare Prediction from Magnetograms: A Controlled Study of Input Design in Deep Learning}
\titlerunning{Solar Flare Prediction from Magnetograms}

\author{ } % Danil Garmaev, Jean Meunier}
\authorrunning{} %D . Garmaev and J. Meunier}

\institute{ } %Universit\'e de Montr\'eal}

\maketitle

\begin{abstract}
% Solar flare forecasting from magnetogram data has been studied with both classical machine learning and deep learning methods. In this work we revisit the problem with a practical objective: to isolate which input-design choices matter most when strong pretrained models are fine-tuned under a common protocol. Using the reduced-resolution active-region dataset of Boucheron et al.~\cite{ref_boucheron}, we study four questions that are useful when preparing inputs for deep-learning models. First, what is the lowest spatial resolution that can be used without a large loss of performance? Second, how much does solar-disk position affect prediction quality? Third, for a 16-frame spatio-temporal model, what temporal cadence provides the most useful context? Fourth, for a fixed temporal span, how many frames are needed? Our main metric is the True Skill Statistic (TSS), with the decision threshold selected on the validation set, and we also report threshold-free ranking metrics such as AUC. The image-based and spatio-temporal experiments were run under a common fine-tuning protocol in order to isolate the effect of the input conditions. The sequence study used an MViT model and was reported for both $\geq C$ and $\geq M$ flare intensity prediction in the next 24 hours. Data augmentation was used as the main training setting. With the completed augmented cadence sweep, the 180-minute setting gave the best mean TSS and AUC for both $\geq C$ and $\geq M$. A subsequent frame-number experiment showed that shorter sequences spanning the same 45-hour window remained competitive, although the preferred number of frames depended on the task and metric. Overall, the results show that temporal context is an important input factor and provide practical guidance for designing efficient and robust deep-learning pipelines for solar flare prediction.
Solar flare forecasting from magnetogram data has been studied with both classical machine learning and deep learning methods. Here we use the reduced-resolution active-region dataset of Boucheron et al.~\cite{ref_boucheron} to study four input-design choices under a common fine-tuning protocol: image resolution, solar-disk position, temporal cadence, and number of frames. The main metric is the True Skill Statistic (TSS), with the decision threshold selected on the validation set; we also report threshold-free ranking metrics such as AUC. Image experiments use a 2D MViT backbone, while sequence experiments use an MViT model for both $\geq C$ and $\geq M$ flare prediction in the next 24 hours. Full resolution performed best on average, but lower resolutions retained much of the performance. Central longitudes gave higher TSS and AUC than outer longitudes, and central-only training improved central-subset TSS. For 16-frame sequences, the 180-minute cadence gave the best mean TSS and AUC for both tasks. When the temporal span was fixed at 45 hours, shorter 4- or 8-frame sequences remained competitive, although the preferred frame count depended on the task and metric. Overall, temporal span was the clearest temporal driver in our experiments, and several input-efficient configurations retained competitive performance.




\keywords{Solar flare prediction \and Magnetogram \and Deep learning \and Vision transformer \and Spatio-temporal learning}
\end{abstract}

\section{Introduction}

The Sun is crucial for life on Earth, being the main source of energy and light. At first sight it looks stable and quiet, a prerequisite to sustain life. However, the Sun's energy output does fluctuate with time, from short periods of hours to much longer periods. This space weather is monitored closely because intense solar activity can affect satellites, telecommunications, navigation systems, and electrical power networks. For this reason, forecasting solar flares is an important problem.

Solar flares are usually produced by magnetic field instabilities associated with active regions (ARs) on the solar surface. While sunspots are visible in optical images, magnetograms give a more direct view of the magnetic field structure; see Fig.~\ref{fig:intro_magnetogram}. This makes magnetograms a natural source of information for automated flare prediction.

\begin{figure}[t]
\centering
\begin{minipage}{0.47\textwidth}
\centering
% \includegraphics[width=\linewidth]{figures/sunspots_1024_20111106.jpg}
\includegraphics[width=\linewidth]{figures/latest.jpg}
\end{minipage}
\hfill
\begin{minipage}{0.47\textwidth}
\centering
\includegraphics[width=\linewidth]{figures/1339_hmi.M_720s.20111106_053600_TAI.1.magnetogram_224.png}
\end{minipage}
\caption{(Left) Full disk magnetogram from the SOHO project website~\cite{ref_soho}. (Right) Magnified view of an active region of the Sun from the Boucheron dataset~\cite{ref_boucheron}. The magnetogram image shows the magnetic field in the solar photosphere, with black and white indicating opposite polarities.}
\label{fig:intro_magnetogram}
\end{figure}


% Machine learning offers many tools for this task. Earlier studies often relied on handcrafted magnetic features and classical classifiers. More recent work applies deep networks directly to AR images or to temporal sequences of AR observations. These approaches include convolutional neural networks, vision transformers, and spatio-temporal architectures. The recent dataset of Boucheron et al.~\cite{ref_boucheron} is especially useful for this line of work because it provides a large and standardized benchmark for image-based flare prediction.

% From a related-work perspective, solar flare forecasting has gradually moved from classical approaches based on handcrafted magnetic features toward deep-learning models operating directly on magnetograms and related solar data~\cite{ref_li_flares,ref_shao_review}. At the same time, benchmarking studies have emphasized standardized evaluation and skill-based metrics such as TSS, since class imbalance can make simpler measures misleading~\cite{ref_bloomfield}. Recent datasets have also made it easier to compare image-based and temporal forecasting approaches under a common protocol~\cite{ref_boucheron}. In this context, our goal is not to introduce another forecasting architecture, but to use established deep models as tools to study how input resolution, disk position, and temporal cadence affect performance.

% From a related-work perspective, the literature has moved from handcrafted magnetic features and classical machine-learning models toward deep models operating directly on AR images or temporal solar observations. Recent work includes convolutional and hybrid deep-learning approaches, such as the fusion model of Li et al.~\cite{ref_li_flares}, while broader reviews describe the growing use of CNNs, recurrent models, vision transformers, and spatio-temporal architectures for flare prediction~\cite{ref_shao_review}. At the same time, benchmarking studies have emphasized standardized evaluation and skill-based metrics such as TSS, since class imbalance can make simpler measures misleading~\cite{ref_bloomfield}. Recent datasets have also made it easier to compare image-based and temporal forecasting approaches under a common protocol~\cite{ref_boucheron}. This context motivates a controlled study of input design choices rather than another architecture comparison.

Machine learning offers many tools for solar flare prediction. Earlier studies often relied on handcrafted magnetic features and classical classifiers, whereas more recent work applies deep models directly to AR images or temporal solar observations. These approaches include convolutional and hybrid deep-learning models, such as the fusion model of Li et al.~\cite{ref_li_flares}, while broader reviews describe the growing use of CNNs, recurrent models, vision transformers, and spatio-temporal architectures for flare prediction~\cite{ref_shao_review}. Benchmarking studies have also emphasized standardized evaluation and skill-based metrics such as TSS, since class imbalance can make simpler measures misleading~\cite{ref_bloomfield}. In this context, the dataset of Boucheron et al.~\cite{ref_boucheron} is especially useful because it provides a large and standardized benchmark, making it possible to study input design choices under a common protocol rather than focusing only on architecture comparisons.


% This paper provides a controlled study of how spatial resolution, solar disk position, temporal cadence, and sequence length affect solar flare prediction when strong pretrained architectures are fine-tuned under a common training protocol. By keeping the model family and optimization setup as consistent as possible, we isolate these input factors and identify configurations that preserve performance while reducing spatial or temporal input cost. 

This paper provides a controlled study of how spatial resolution, solar disk position, temporal cadence, and sequence length affect solar flare prediction when strong pretrained architectures are fine-tuned under a common training protocol. By isolating these input factors, we identify practical configurations that preserve performance while reducing spatial or temporal input cost. Our results show that temporal span is more important than dense sampling in the sequence experiments, while lower spatial resolution and sparse temporal inputs can remain competitive. We focus on four questions. First, what is the lowest AR image resolution that can be used without a significant loss of performance? Second, does the position of the AR on the solar disk, in longitude and latitude, have a measurable impact on performance? Third, for a fixed 16-frame sequence model, what temporal cadence provides the most useful temporal context? Fourth, for a fixed temporal span, how many frames are needed for effective prediction?

\section{Dataset}

For this work we use the reduced-resolution solar flare prediction dataset proposed by Boucheron et al.~\cite{ref_boucheron}. According to Boucheron et al., the reduced-resolution version retains comparable classification performance to the full-resolu-tion dataset~\cite{ref_boucheron}. Line-of-sight magnetograms of solar active regions are provided by the Helioseismic and Magnetic Imager (HMI) on board NASA's Solar Dynamics Observatory. The dataset covers the period from May 1, 2010 to December 31, 2018 and includes labels for flare occurrence in a 24-hour prediction window.

We use the preconfigured $224 \times 224$ AR images. The subset is restricted to ARs whose centers are within $\pm 60^\circ$ latitude and longitude, which reduces projection effects and removes the most problematic edge cases. The original data split by active region is kept unchanged, with separate training, validation, and testing sets. According to the dataset description, the reduced-resolution subset contains 759{,}357 training images, 95{,}933 validation images, and 94{,}757 test images.

Two binary prediction tasks were studied. In the $\geq C$ task, all flares of class C or higher are considered positive. In the training split, this corresponds to 149{,}249 positive images out of 759{,}357, or 19.7\%. In the $\geq M$ task, only more intense flares of class M or higher are positive; this corresponds to 23{,}395 positive training images, or 3.1\%. Class M therefore requires capturing more complex and less frequent processes. Consequently, the second task is more imbalanced and more difficult.

\section{Performance Metric}

In solar flare forecasting the data are imbalanced, so accuracy alone is not sufficient to evaluate a model. In this paper we use the True Skill Statistic (TSS), which is widely recommended for flare prediction studies~\cite{ref_bloomfield}. It is defined from the confusion matrix as
\begin{equation}
\mathrm{TSS} = \frac{TP}{TP + FN} - \frac{FP}{FP + TN} = \text{Sensitivity} + \text{Specificity} -1 .
\end{equation}

The TSS ranges from $-1$ to $1$, with $1$ being the best value. It can also be written as sensitivity plus specificity minus one.

In our experiments the main metric is TSS. More precisely, throughout the paper the reported TSS is obtained with the decision threshold selected on the validation set, subject to minimum recall (0.5) and precision (0.15) constraints to avoid degenerate solutions. 

We also report the area under the ROC curve (AUC). Unlike TSS, AUC does not depend on a single operating threshold and is useful for measuring how well the model separates positive and negative samples over all possible thresholds. The area under the precision-recall curve (PR-AUC) is included because it is informative for the more imbalanced $\geq M$ task. Finally, we report F1 as a compact threshold-dependent summary of precision and recall.

\section{Method}

All experiments were designed to isolate input-related factors while keeping the architecture family and the fine-tuning protocol as consistent as possible.

\subsection{Model Choice and Multiscale Attention}

The model choice was motivated by the need to capture both local magnetic structures and global active-region organization. Multiscale Vision Transformers (MViT) build hierarchical feature pyramids by reducing token resolution while increasing channel dimensionality~\cite{ref_mvit}. This is relevant for magnetograms, where flare-related signals may arise from both fine-scale polarity inversion structures and larger-scale magnetic configurations.

MViT-family models can be applied to both image and video recognition. In the image experiments, we used a 2D MViTv2-style backbone selected after preliminary screening of single-image architectures; MViTv2 extends MViT with decomposed relative positional embeddings and residual pooling connections~\cite{ref_mvitv2}. In the sequence experiments, we used a video MViT backbone from the available training pipeline, which directly supports spatio-temporal tokenization and the variable temporal lengths required for the frame-count study.

We do not treat MViT versus MViTv2 as an architecture comparison. Instead, both are used as closely related multiscale-attention backbones to study input design: spatial resolution, disk position, temporal cadence, and sequence length.


\subsection{Data Augmentation}

Following common practice in recent deep-learning studies, we used augmented training as the main setting in both the image-based and sequence-based experiments. The augmentation pipeline applied random horizontal and vertical flips, random rotations up to $\pm 30^\circ$, magnetic polarity inversion, and small integer noise in the range $[-5,5]$ (after magnetogram intensities were normalized to $[0,255]$). Random rotations were limited to avoid unrealistic data generation since ARs usually appear as pairs of opposite magnetic polarities aligned east-west (a $\pm90^\circ$ rotation would produce a north-south orientation). When a spatial-resolution ablation was used, the downsampling step was applied before augmentation so that the transforms did not reintroduce fine details that had already been removed.

For sequence models, the same geometric transform and the same polarity inversion were applied consistently to all frames of a sequence, while the small integer noise could vary from frame to frame. This kept the temporal structure coherent while still perturbing the appearance of the input.



\subsection{Resolution for AR Images}

The first experiment studied the effect of spatial resolution for image-based flare prediction. Before fixing this experiment, we screened several single-image backbones, including ConvNeXt-Tiny~\cite{ref_convnext}, Swin-Tiny~\cite{ref_swin}, ViT-Small/16~\cite{ref_vit}, and a 2D MViT variant from the same model family as the sequence experiment~\cite{ref_mvit}. We then retained the 2D MViT model for the image-based studies in order to keep the controlled comparisons within a common architectural family.

The original AR images are available at $224 \times 224$. To study lower resolutions, we downsample them to $112 \times 112$ and $56 \times 56$, then upsample them back to $224 \times 224$ before sending them to the backbone. This keeps the network architecture unchanged while allowing us to control how much spatial detail is removed from the magnetogram. The goal is to estimate the lowest useful resolution for flare prediction.

Training followed the shared setup described in Sect.~\ref{sec:repro}, and the main comparison metric was the True Skill Statistic (TSS).

\subsection{Longitude and Latitude on the Solar Disk}

The second experiment studied the effect of AR position on the solar disk. Because the Sun is a sphere, active regions far from disk center are more affected by projection and geometric distortion. Our hypothesis was that prediction should be easier for ARs located closer to the center of the disk, for example within $\pm 30^\circ$, than for ARs located farther away.

This experiment was built on top of the first one, since it reused the image model and the resolution selected there. It was organized around two related but distinct comparisons. First, we evaluated the same model on different longitude ranges, in particular $|$longitude$| \leq 30^\circ$ and $30^\circ < |$longitude$| \leq 60^\circ$, in order to measure how performance changed with disk position under a fixed training setup. Second, we compared a model trained only on the central $|$longitude$| \leq 30^\circ$ subset against the baseline model trained on the full $\pm 60^\circ$ set, with both models evaluated on the same central-region test subset. This second comparison isolates the effect of training specialization and should therefore be interpreted separately from the longitude-bucket comparison. Seed variation was used to quantify uncertainty in both views.

\subsection{Time Interval Between Frames for AR Sequences}

For the sequence experiments we used a 16-frame MViT model~\cite{ref_mvit}. Each input is a temporal sequence of 16 magnetograms of the same active region, still at $224 \times 224$ resolution. Training again followed the shared setup described in Sect.~\ref{sec:repro}.

We study three cadences: 60, 120, and 180 minutes between frames. With 16 frames, these correspond to temporal contexts of 15, 30, and 45 hours, respectively. The goal is to see whether longer temporal context helps the model or whether a shorter cadence is sufficient.

Every sequence setting was repeated with five random seeds. As in the single-image study, the main comparison metric is TSS.

\subsection{Number of Frames for AR Sequences}

The final experiment studied the number of frames needed when the temporal span was kept fixed. We used the 45-hour span of the 16-frame, 180-minute-cadence setting because this was the strongest setting in the cadence study. Unlike the initial fixed-input design, the MViT implementation used here can accept different temporal lengths, so no interpolation, padding, or repetition of frames was required. We therefore evaluated sequences with $T=4$, $T=8$, and $T=16$ frames directly by selecting evenly spaced temporal subsets that covered the same 45-hour window. For example, the $T=4$ setting used frames at 45, 30, 15, and 0 hours before the prediction time.

This experiment separates two related aspects of temporal information: the total temporal span covered by the sequence and the number of observations provided to the model within that span. It tests whether the dense 16-frame sequence is necessary, or whether a sparser sequence can preserve most of the predictive information while reducing the temporal input size.






\subsection{Reproducibility and Training Setup}
\label{sec:repro}

To keep the study reproducible while limiting the number of changing factors, we used a common training recipe across the main experiments. Training followed the same general optimization strategy for both the image and sequence models. We used supervised fine-tuning of pretrained models with AdamW, weight decay of 0.01, weighted binary cross-entropy, and a cosine learning-rate schedule with differential learning rates: $10^{-6}$ for the pretrained backbone and $10^{-4}$ for the classification head. The image models used a batch size of 32 and a 2-epoch warmup, while the sequence models used a batch size of 8 and a 3-epoch warmup. The maximum training budget was 50 epochs, and model selection was based on validation TSS with early stopping to limit overfitting. The augmentation pipeline combined horizontal and vertical flips, random rotations up to $\pm 30^\circ$, magnetic polarity inversion, and small integer noise. For sequence inputs, the geometric transforms and polarity inversion were shared across all frames of the same sequence in order to preserve temporal coherence. The reported test-set TSS values were computed with the decision threshold selected on the validation set. Table~\ref{tab:hparams} summarizes the main settings used for the selected image and sequence models.


% To keep the study reproducible while limiting the number of changing factors, we used a common training recipe across the main experiments. For the image experiments, we first screened several single-image backbones, including ConvNeXt-Tiny~\cite{ref_convnext}, Swin-Tiny~\cite{ref_swin}, ViT-Small/16~\cite{ref_vit}, and a 2D MViT variant from the same model family as the sequence model~\cite{ref_mvit}. We then retained the 2D MViT model for the image-based studies. For the sequence experiments, we used a pretrained 16-frame MViT model. In all cases, the backbone input remained $224 \times 224$. For the resolution study, the lower-resolution variants were obtained by downsampling the AR image to $112 \times 112$ or $56 \times 56$ and then upsampling it back to $224 \times 224$ before the backbone; this downsampling step was applied before augmentation.

% Training followed the same general optimization strategy for both the image and sequence models. We fine-tuned pretrained models with AdamW, weight decay of 0.01, weighted binary cross-entropy, and a cosine learning-rate schedule with differential learning rates: $10^{-6}$ for the pretrained backbone and $10^{-4}$ for the classification head. The image models used a batch size of 32 and a 2-epoch warmup, while the sequence models used a batch size of 8 and a 3-epoch warmup. The maximum training budget was 50 epochs, and model selection was based on validation TSS with early stopping to limit overfitting. The augmentation pipeline combined horizontal and vertical flips, random rotations up to $\pm 30^\circ$, magnetic polarity inversion, and small integer noise. For sequence inputs, the geometric transforms and polarity inversion were shared across all frames of the same sequence in order to preserve temporal coherence. The reported test-set TSS values were computed with the decision threshold selected on the validation set. Table~\ref{tab:hparams} summarizes the main settings used for the selected image and sequence models.


\begin{table}[H]
\caption{Main training settings used for the selected image and sequence models.}
\label{tab:hparams}
\centering
\scriptsize
\renewcommand{\arraystretch}{0.95}
\setlength{\tabcolsep}{2.5pt}
\begin{tabular}{@{}>{\raggedright\arraybackslash}p{0.20\textwidth}>{\raggedright\arraybackslash}p{0.34\textwidth}>{\raggedright\arraybackslash}p{0.34\textwidth}@{}}
\toprule
Item & Image model & Sequence model \\
\midrule
Model & 2D MViT (pretrained) & MViT (pretrained) \\
Input & $1 \times 224 \times 224$ & $T \times 1 \times 224 \times 224$ \\
% Resolution / cadence & $224$ input; for the resolution study, downsample to $112$ or $56$, then upsample back to $224$ & 16 frames; 60, 120, or 180 min cadence \\
Augmentation & H/V flips, rotation up to $\pm 30^\circ$, polarity inversion, integer noise & Same transforms; geometric operations shared across frames \\
Optimizer & AdamW & AdamW \\
Backbone LR & $10^{-6}$ & $10^{-6}$ \\
Head LR & $10^{-4}$ & $10^{-4}$ \\
Weight decay & 0.01 & 0.01 \\
Scheduler & Cosine & Cosine \\
Warmup & 2 epochs & 3 epochs \\
Loss & Weighted BCE & Weighted BCE \\
Batch size & 32 & 8 \\
Epochs & 50 & 50 \\
\bottomrule
\end{tabular}
\end{table}


\section{Results}

In the result tables, bold values indicate the best mean value for each metric within a given flare-intensity task.

\subsection{Resolution for AR Images}

Table~\ref{tab:single} presents the resolution study. It compares the same image model at $224 \times 224$, $112 \times 112$, and $56 \times 56$. The goal is to see whether reducing spatial detail causes only a small drop in TSS or whether there is a clear performance break at low resolution.

\begin{table}[H]
\caption{Results for the A1 resolution study with the selected 2D MViT image model.}
\label{tab:single}
\centering
\scriptsize
\setlength{\tabcolsep}{5pt}
\begin{tabular}{@{}ccccc@{}}
\toprule
Resolution & TSS & AUC & PR-AUC & F1 \\
\midrule
\multicolumn{5}{c}{$\geq C$ task} \\
\midrule
224 & $\best{0.535 \pm 0.024}$ & $\best{0.848 \pm 0.012}$ & $\best{0.670 \pm 0.011}$ & $\best{0.606 \pm 0.010}$ \\
112 & $0.502 \pm 0.057$ & $0.831 \pm 0.032$ & $0.646 \pm 0.030$ & $0.585 \pm 0.033$ \\
56 & $0.510 \pm 0.054$ & $0.831 \pm 0.030$ & $0.625 \pm 0.049$ & $0.582 \pm 0.033$ \\
\addlinespace
\multicolumn{5}{c}{$\geq M$ task} \\
\midrule
224 & $\best{0.641 \pm 0.020}$ & $\best{0.911 \pm 0.008}$ & $0.378 \pm 0.062$ & $0.268 \pm 0.013$ \\
112 & $0.594 \pm 0.040$ & $0.890 \pm 0.012$ & $0.384 \pm 0.021$ & $\best{0.269 \pm 0.047}$ \\
56 & $0.580 \pm 0.012$ & $0.885 \pm 0.006$ & $\best{0.395 \pm 0.036}$ & $0.264 \pm 0.018$ \\
\bottomrule
\end{tabular}
\end{table}


The performance drop from full resolution to lower resolution was moderate rather than abrupt, which suggests that part of the flare-related information is still preserved at coarse scales. For the $\geq C$ task, the $112 \times 112$ and $56 \times 56$ variants remained relatively close to full resolution in TSS and AUC. For the harder $\geq M$ task, however, the full-resolution input gave the best mean TSS and AUC among the tested settings. The PR-AUC values should be read relative to class prevalence: the $\geq C$ task has about 20\% positives, whereas the $\geq M$ task has only about 3\% positives, so PR-AUC values around 0.37-0.39 for $\geq M$ are far above the no-skill baseline even though they are numerically lower than the $\geq C$ values.


\subsection{Longitude and Latitude on the Solar Disk}

The second experiment studied the effect of position on the solar disk. We expected ARs within about $\pm 30^\circ$ of disk center to give better performance than ARs located farther away, because geometric distortion is weaker there. Pandey et al.~\cite{ref_long} also reported a decline in model performance with increasing longitudinal coverage. We focused on longitude only since most ARs are already within $\pm30^\circ$ latitude. Table~\ref{tab:a2} reports the baseline model evaluated on the full $\pm60^\circ$ test set, the same baseline model evaluated on two disjoint longitude buckets, and the model trained only on the central $|$longitude$| \leq 30^\circ$ subset and evaluated on that same central test bucket. Fig.~\ref{fig:a2_graph} summarizes the two main TSS comparisons: bucketed evaluation for the baseline model and central-only training on the central test subset.


\begin{table}[H]
\caption{Results for the A2 longitude experiment.}
\label{tab:a2}
\centering
\scriptsize
\setlength{\tabcolsep}{3pt}
\resizebox{\textwidth}{!}{%
\begin{tabular}{@{}llcccc@{}}
\toprule
Train longitudes & Test longitudes & TSS & AUC & PR-AUC & F1 \\
\midrule
\multicolumn{6}{c}{$\geq C$ task} \\
\midrule
Full $\pm60^\circ$ & Full $\pm60^\circ$ & $0.535 \pm 0.024$ & $0.848 \pm 0.012$ & $0.670 \pm 0.011$ & $0.606 \pm 0.010$ \\
Full $\pm60^\circ$ & $|$lon$| \leq 30^\circ$ & $0.545 \pm 0.023$ & $\best{0.859 \pm 0.006}$ & $0.637 \pm 0.013$ & $0.591 \pm 0.007$ \\
Full $\pm60^\circ$ & $30^\circ < |$lon$| \leq 60^\circ$ & $0.506 \pm 0.031$ & $0.832 \pm 0.021$ & $\best{0.707 \pm 0.028}$ & $\best{0.622 \pm 0.018}$ \\
$|$lon$| \leq 30^\circ$ & $|$lon$| \leq 30^\circ$ & $\best{0.566 \pm 0.029}$ & $\best{0.859 \pm 0.015}$ & $0.601 \pm 0.051$ & $0.590 \pm 0.025$ \\
\addlinespace
\multicolumn{6}{c}{$\geq M$ task} \\
\midrule
Full $\pm60^\circ$ & Full $\pm60^\circ$ & $0.641 \pm 0.020$ & $0.911 \pm 0.008$ & $0.378 \pm 0.062$ & $\best{0.268 \pm 0.013}$ \\
Full $\pm60^\circ$ & $|$lon$| \leq 30^\circ$ & $0.580 \pm 0.125$ & $\best{0.927 \pm 0.012}$ & $\best{0.417 \pm 0.077}$ & $0.242 \pm 0.086$ \\
Full $\pm60^\circ$ & $30^\circ < |$lon$| \leq 60^\circ$ & $0.571 \pm 0.148$ & $0.894 \pm 0.016$ & $0.353 \pm 0.024$ & $0.228 \pm 0.067$ \\
$|$lon$| \leq 30^\circ$ & $|$lon$| \leq 30^\circ$ & $\best{0.688 \pm 0.039}$ & $0.923 \pm 0.015$ & $0.337 \pm 0.079$ & $0.245 \pm 0.014$ \\
\bottomrule
\end{tabular}
}
\end{table}


\begin{figure}[H]
\centering
\includegraphics[width=\textwidth]{figures/a2_longitude_tss_std.png}
\caption{A2 longitude experiment summarized by TSS. Panel A compares the baseline model across inner and outer longitude buckets. Panel B compares full-disk and central-only training on the same central test subset. Error bars show standard deviation across five seeds.}
\label{fig:a2_graph}
\end{figure}

For both label settings, the baseline model gave higher mean TSS and AUC on the inner longitude bucket than on the outer bucket, which is consistent with weaker projection effects near disk center. The full-test-set row is included as an overall reference, while the two disjoint bucket rows isolate how the same trained model behaves as a function of longitude.


When we restricted training to the central $|$longitude$| \leq 30^\circ$ subset and evaluated on the same central test bucket, the mean TSS increased from 0.545 to 0.566 for $\geq C$ and from 0.580 to 0.688 for $\geq M$. This suggests that training specialization can further help on central-region data, especially for the harder $\geq M$ task, although the effect is metric-dependent: for example, the baseline model still gave the highest mean AUC and PR-AUC for the $\geq M$ central bucket.




\subsection{Cadence Study with 16-Frame MViT}

Table~\ref{tab:video} summarizes the completed A3 sequence experiments with augmented training. For the $\geq C$ task, the mean TSS increased monotonically with cadence, from 0.664 at 60 minutes to 0.745 at 120 minutes and 0.789 at 180 minutes. This indicates that, within the tested range, temporal span mattered more than dense short-interval sampling.


For the $\geq M$ task, the 180-minute cadence also gave the highest mean TSS and the highest mean AUC, although the TSS margin over 120 minutes was smaller and the variance remained larger for this harder task. Taken together, the completed augmented runs favor the longest tested temporal context. With 16 frames, the 180-minute setting spans 45 hours of active-region evolution, compared with 30 hours at 120 minutes and 15 hours at 60 minutes, which suggests that predictive signal is carried more by multi-hour magnetic evolution than by near-duplicate short-cadence observations.


\begin{table}[H]
\caption{A3 cadence study for the 16-frame MViT model.}
\centering
\label{tab:video}
\scriptsize
\setlength{\tabcolsep}{4pt}
\begin{tabular}{@{}ccccc@{}}
\toprule
Cadence (min) & TSS & AUC & PR-AUC & F1 \\
\midrule
\multicolumn{5}{c}{$\geq C$ task} \\
\midrule
60  & $0.664 \pm 0.025$ & $0.904 \pm 0.020$ & $0.748 \pm 0.039$ & $0.664 \pm 0.017$ \\
120 & $0.745 \pm 0.032$ & $0.935 \pm 0.014$ & $0.804 \pm 0.044$ & $0.724 \pm 0.016$ \\
180 & $\best{0.789 \pm 0.020}$ & $\best{0.957 \pm 0.004}$ & $\best{0.842 \pm 0.027}$ & $\best{0.740 \pm 0.014}$ \\
\addlinespace
\multicolumn{5}{c}{$\geq M$ task} \\
\midrule
60  & $0.644 \pm 0.062$ & $0.900 \pm 0.012$ & $0.398 \pm 0.091$ & $0.295 \pm 0.049$ \\
120 & $0.671 \pm 0.052$ & $0.914 \pm 0.007$ & $\best{0.399 \pm 0.055}$ & $0.314 \pm 0.035$ \\
180 & $\best{0.674 \pm 0.134}$ & $\best{0.947 \pm 0.011}$ & $0.393 \pm 0.042$ & $\best{0.377 \pm 0.090}$ \\
\bottomrule
\end{tabular}
\end{table}

The PR-AUC contrast between the two label settings mainly reflects their different positive rates. For $\geq C$, the no-skill PR-AUC baseline is about 0.20, so values near 0.85 indicate strong ranking performance. For $\geq M$, the baseline is about 0.03, making the observed PR-AUC values near 0.39 roughly an order of magnitude above random performance despite the lower absolute value.

\subsection{Number of Frames over a Fixed 45-Hour Window}

Table~\ref{tab:frames} summarizes the frame-number experiment. The temporal span was fixed at 45 hours, matching the 16-frame, 180-minute-cadence setting, and the MViT model was evaluated with $T=4$, $T=8$, and $T=16$ input frames sampled from this window. Since TSS and F1 depend on a selected operating threshold, PR-AUC is included here together with ROC-AUC to provide threshold-free views of performance, with PR-AUC being especially relevant for the imbalanced $\geq M$ task and its much lower positive-rate baseline.

For the $\geq C$ task, the sparser sequences performed at least as well as the 16-frame setting. The best mean TSS and ROC-AUC were obtained with $T=4$, while the best mean PR-AUC was obtained with $T=8$; however, the AUC values were very close across all three sequence lengths. For the harder $\geq M$ task, $T=8$ gave the highest mean TSS and ROC-AUC, whereas $T=16$ retained the highest mean PR-AUC and a nearly identical F1 score. These results identify a practical input-efficient regime: the full 16-frame input was not always necessary when the 45-hour span was preserved, although no single frame count was best across all metrics.

\begin{table}[H]
% \caption{Results for the frame-number study over a fixed 45-hour window with the MViT sequence model.}
\caption{Effect of frame count over a fixed 45-hour window (MViT).}
\centering
\label{tab:frames}
\scriptsize
\setlength{\tabcolsep}{6pt}
\begin{tabular}{@{}ccccc@{}}
\toprule
Frames & TSS & AUC & PR-AUC & F1 \\
\midrule
\multicolumn{5}{c}{$\geq C$ task} \\
\midrule
16 & $0.789 \pm 0.020$ & $0.957 \pm 0.004$ & $0.842 \pm 0.027$ & $0.740 \pm 0.014$ \\
8  & $0.800 \pm 0.044$ & $0.960 \pm 0.008$ & $\best{0.865 \pm 0.022}$ & $0.732 \pm 0.038$ \\
4  & $\best{0.810 \pm 0.032}$ & $\best{0.961 \pm 0.007}$ & $0.861 \pm 0.019$ & $\best{0.741 \pm 0.032}$ \\
\addlinespace
\multicolumn{5}{c}{$\geq M$ task} \\
\midrule
16 & $0.674 \pm 0.134$ & $0.947 \pm 0.011$ & $\best{0.393 \pm 0.042}$ & $\best{0.377 \pm 0.090}$ \\
8  & $\best{0.758 \pm 0.102}$ & $\best{0.952 \pm 0.010}$ & $0.364 \pm 0.076$ & $0.376 \pm 0.048$ \\
4  & $0.752 \pm 0.075$ & $0.946 \pm 0.012$ & $0.361 \pm 0.071$ & $0.330 \pm 0.051$ \\
\bottomrule
\end{tabular}
\end{table}

% \section{Discussion}

% This study complements architecture-focused work by isolating input-design choices under a controlled fine-tuning setup. Four findings stand out. First, temporal context was the strongest factor among the tested sequence settings: for 16-frame sequences, the 180-minute cadence gave the best mean TSS and AUC for both labels. Second, with the temporal span fixed at 45 hours, sparser 4- or 8-frame sequences remained competitive, suggesting that temporal coverage can matter more than the number of sampled frames. Third, reducing spatial resolution caused only moderate loss, although full resolution remained best on average, especially for $\geq M$. Fourth, longitude mattered: central regions gave higher mean TSS and AUC, and central-only training improved central-subset TSS.

% The cadence results suggest that temporal span, rather than sampling density alone, was the main driver of the sequence-model gains in this setting. With 16 frames, the 180-minute cadence spans 45 hours of active-region evolution, compared with 30 hours at 120 minutes and 15 hours at 60 minutes. This interpretation is clearest for $\geq C$, where mean TSS increased steadily from 60 to 120 to 180 minutes. For $\geq M$, the best mean TSS and AUC were also obtained at 180 minutes, although the improvement over 120 minutes was smaller and accompanied by larger variance.

% The frame-number experiment adds an important nuance: preserving the temporal span was often more important than keeping all 16 observations. For $\geq C$, 4 frames gave the highest mean TSS and ROC-AUC, while 8 frames gave the highest mean PR-AUC. For $\geq M$, 8 frames gave the highest mean TSS and ROC-AUC, whereas 16 frames retained the highest mean PR-AUC. Since the AUC differences were small and the preferred setting changed with the metric, these results should be interpreted as evidence that sparse sampling over a long temporal window is viable, not as a definitive ranking of frame counts. This also avoids artificial interpolation or frame repetition.

% The comparison with the native-resolution single-image baseline further supports the importance of temporal information. For $\geq C$, the best spatio-temporal setting improved TSS from 0.535 to 0.789 and AUC from 0.848 to 0.957. For the harder $\geq M$ task, the TSS gain was smaller, from 0.641 to 0.674, but AUC still increased from 0.911 to 0.947. These results are consistent with the idea that flare-relevant magnetic evolution often unfolds over several hours rather than only through rapid frame-to-frame fluctuations.

% PR-AUC values must be interpreted relative to the class prevalence. For the $\geq C$ task, the positive rate is approximately 20\%, leading to comparatively higher PR-AUC values. In contrast, the $\geq M$ task has a positive rate of about 3\%, which defines a substantially lower no-skill baseline. Consequently, PR-AUC values in the range of 0.3--0.4 for the $\geq M$ task represent a significant improvement over the baseline, corresponding to an order-of-magnitude gain relative to random performance, despite being numerically lower than those observed for the $\geq C$ task.

% The image-based experiments give complementary pipeline guidance. The resolution study suggests that large-scale magnetic organization carries much of the predictive information, even if finer detail still matters, particularly for $\geq M$. The longitude study is consistent with projection effects making learned features less reliable away from disk center. Such effects can in principle be reduced through 3D reprojection methods~\cite{ref_brown_walker}, but this adds preprocessing and interpolation. The two longitude comparisons should also be interpreted separately: bucketed evaluation measures disk-position sensitivity of the same model, while central-only training tests specialization to central-region data. In our results, central-only training improved central-subset TSS, especially for $\geq M$, although the effect remained metric-dependent.

% The main limitations are that the image-based studies are more compact than the cadence analysis and that the experiments were not intended to exhaustively optimize architectures or training details. Future work should test additional cadence and window combinations, extend the longitude-specialization analysis, combine magnetograms with handcrafted magnetic features, and examine whether filtering overlapping ARs improves label reliability. Longer temporal windows were not included because they require each AR to remain observable with enough valid frames across the full window, which reduces the number of usable training sequences. This reduction is important for deep sequence models, which need enough independent examples to avoid overfitting.


\section{Discussion}

This study complements architecture-focused work by isolating input-design choices under a controlled fine-tuning setup. Four findings stand out: the 180-minute cadence gave the best mean TSS and AUC for both labels; with the temporal span fixed at 45 hours, 4- or 8-frame sequences remained competitive; reducing spatial resolution caused only moderate loss, although full resolution remained best on average, especially for $\geq M$; and longitude mattered, with central regions giving higher mean TSS and AUC.

The cadence results suggest that temporal span, rather than sampling density alone, was the main driver of the sequence-model gains. With 16 frames, the 180-minute cadence spans 45 hours of active-region evolution, compared with 30 hours at 120 minutes and 15 hours at 60 minutes. This interpretation is clearest for $\geq C$, where mean TSS increased steadily with cadence; for $\geq M$, the best mean TSS and AUC were also obtained at 180 minutes, although with larger variance.

The frame-number experiment adds an important nuance: preserving temporal span was often more important than keeping all 16 observations. Since the AUC differences were small and the preferred setting changed with the metric, these results should be interpreted as evidence that sparse sampling over a long temporal window is viable, not as a definitive ranking of frame counts.

The comparison with the native-resolution single-image baseline further supports the importance of temporal information. For $\geq C$, the best spatio-temporal setting improved TSS from 0.535 to 0.789 and AUC from 0.848 to 0.957. For $\geq M$, the TSS gain was smaller, from 0.641 to 0.674, but AUC still increased from 0.911 to 0.947. These results are consistent with flare-relevant magnetic evolution unfolding over several hours rather than only through rapid frame-to-frame fluctuations.

The image-based experiments give complementary pipeline guidance. The resolution study suggests that large-scale magnetic organization carries much of the predictive information, even if finer detail still matters, particularly for $\geq M$. The longitude study is consistent with projection effects making learned features less reliable away from disk center. The two longitude comparisons should be interpreted separately: bucketed evaluation measures disk-position sensitivity of the same model, while central-only training tests specialization to central-region data. In our results, central-only training improved central-subset TSS, especially for $\geq M$, although the effect remained metric-dependent.

The main limitations are that the image-based studies are more compact than the cadence analysis and that the experiments were not intended to exhaustively optimize architectures or training details. Future work should test additional cadence and window combinations, extend longitude-specialization analysis, combine magnetograms with handcrafted magnetic features, and examine whether filtering overlapping ARs improves label reliability. Longer temporal windows were not included because they require each AR to remain observable with enough valid frames, reducing the number of usable training sequences.



% \section{Discussion}

% This study focuses on understanding input design choices under a controlled deep-learning setting, complementing existing work that emphasizes architectural innovation. Nevertheless, the observed performance range was broadly consistent with recent deep-learning studies and reviews in this area~\cite{ref_li_flares,ref_shao_review}. Four main findings emerge from the experiments. First, temporal context was the strongest input factor among the tested sequence settings: the completed augmented cadence runs favored the 180-minute cadence for both $\geq C$ and $\geq M$ in both TSS and AUC. Second, when the temporal span was fixed at 45 hours, sparser sequences remained competitive, although the best frame count depended on whether threshold-dependent metrics or threshold-free metrics were emphasized. Third, reducing spatial resolution caused only a moderate loss of performance, even though full resolution remained best on average, especially for the harder $\geq M$ task. Fourth, disk position mattered: the longitude bucket closer to disk center gave higher mean TSS and AUC, and central-only training further improved TSS on the central test subset.

% The cadence results indicate that longer temporal coverage is more valuable than denser temporal sampling. With 16 frames, the 180-minute setting spans 45 hours of active-region evolution, compared with 30 hours at 120 minutes and 15 hours at 60 minutes. The fact that 180 minutes performed best therefore suggests that flare prediction in this setting depends more on capturing broad multi-hour magnetic evolution than on resolving highly redundant short-interval changes. This interpretation is especially clear for the $\geq C$ task, where mean TSS increased steadily from 60 to 120 to 180 minutes. For $\geq M$, the best mean TSS and AUC were also obtained at 180 minutes, although the improvement over 120 minutes was smaller and accompanied by larger variance.

% The frame-number experiment adds an important nuance to this conclusion. Here the 45-hour temporal span was preserved while the number of sampled frames was reduced from 16 to 8 or 4. For $\geq C$, 4 frames gave the highest mean TSS and ROC-AUC, while 8 frames gave the highest mean PR-AUC. For $\geq M$, 8 frames gave the highest mean TSS and ROC-AUC, whereas 16 frames retained the highest mean PR-AUC. Since the AUC differences were small and the preferred setting changed with the metric, this experiment should be interpreted as evidence that sparse sampling over a long temporal window is viable, not as a definitive ranking of frame counts. This also suggests that the variable-length MViT can use shorter sequences effectively without artificial interpolation or repetition. Longer sequences spanning several days were not considered, since this would substantially reduce the available training set and would be harder to deploy in practice because of model input-size limits and AR sequences with too few valid images.

% The comparison with the native-resolution single-image baseline further supports the importance of temporal information. For $\geq C$, the best spatio-temporal model substantially outperformed the single-image model, improving TSS from 0.535 to 0.789 and AUC from 0.848 to 0.957. For the more difficult $\geq M$ task, the TSS improvement was smaller, from 0.641 to 0.674, but the AUC still increased clearly, from 0.911 to 0.947. These results are physically plausible, since flare-relevant magnetic energy buildup and structural reorganization often unfold over many hours rather than only through rapid frame-to-frame fluctuations. Whether even longer temporal contexts remain beneficial remains an open question for future work.

% The image-based experiments remain important because they speak directly to pipeline design. In A1, full resolution remained best on average, but the lower-resolution inputs retained a substantial fraction of the performance. This is consistent with the idea that large-scale magnetic organization carries much of the predictive information, even if finer spatial detail still matters, particularly for $\geq M$. In A2, the bucket closer to disk center gave the higher mean TSS and AUC for both label settings, although the outer bucket often retained higher recall, so the longitude effect was not uniform across all operating metrics. This pattern is consistent with projection effects making learned features less reliable away from disk center. Such effects can in principle be reduced through 3D transformations that reproject ARs on a sphere~\cite{ref_brown_walker}, but doing so introduces additional preprocessing steps and interpolation.

% For the longitude study, it is important to distinguish between two questions: whether the same model performs differently across disk-position buckets, and whether a model specialized to central regions improves performance on central-region data. These are both meaningful comparisons, but they answer different scientific questions and should therefore be interpreted separately. In our results, the central-only trained model further improved mean TSS on the central test subset, especially for $\geq M$, although the associated variance remained substantial.



% There are also limitations. The image-based studies remain more compact than the cadence analysis, and their conclusions should therefore be interpreted with some caution. In addition, the experiments were designed to study model behavior under controlled changes of the input rather than to exhaustively optimize every architecture and training detail.

% Several improvements can be explored next. The most direct one is to extend the resolution and position studies with more seeds and a more complete analysis of longitude specialization. Beyond that, it would be useful to test additional combinations of cadence and sequence length, study whether combining magnetogram images with handcrafted magnetic features can improve robustness, and test models trained from scratch on this domain instead of relying only on fine-tuning from natural-image or video datasets. Finally, the dataset itself could be improved. For instance, the dataset contains some overlapping ARs that could introduce errors since the machine learning algorithm learns the flaring behavior of the central AR only. Exclusion of these images could be considered in our future work.


\section{Conclusion}


We presented a controlled study of four practical input-design questions for solar flare prediction from magnetogram data: image resolution, AR position on the solar disk, temporal cadence, and number of frames for spatio-temporal models. Augmentation was used as the main training setting. The A1 results showed that reducing resolution caused only a moderate loss of performance, although full resolution remained best on average, especially for $\geq M$. The cadence results showed that temporal context was consistently useful, with the 180-minute cadence giving the best mean TSS and AUC for both $\geq C$ and $\geq M$ in our augmented runs. The frame-number study then showed that, for the same 45-hour temporal span, sparser sequences with 4 or 8 frames could remain competitive with the 16-frame setting, while the preferred frame count depended on the metric being considered. The A2 analysis further suggested that active regions closer to disk center tend to give better mean TSS and AUC, and that central-only training can further improve TSS on the central test subset. Together, these findings provide practical guidance for designing efficient flare-prediction pipelines and motivate a more complete analysis of spatial detail, disk-position effects, and temporal context.

\begin{credits}
\subsubsection{\discintname}
The authors have no competing interests to declare that are relevant to the content of this article.
\end{credits}

\begin{thebibliography}{11}

\bibitem{ref_boucheron}
Boucheron, L.E., Vincent, T., Grajeda, J.A., Wuest, E.: Solar active region magnetogram image dataset for studies of space weather. Sci. Data \textbf{10}, 825 (2023). \doi{10.1038/s41597-023-02628-8}

\bibitem{ref_bloomfield}
Bloomfield, D.S., Higgins, P.A., Gallagher, P.T., McAteer, R.T.J.: Toward reliable benchmarking of solar flare forecasting methods. Astrophys. J. Lett. \textbf{747}(2), L41 (2012). \doi{10.1088/2041-8205/747/2/L41}

\bibitem{ref_li_flares}
Li, Y.Y., Huang, S.Y., Xu, S.B., Yuan, Z.G., Jiang, K., Xiong, Q.Y., Lin, R.T.: Solar flare forecasting based on a fusion model. Earth Planet. Phys. \textbf{9}(1), 171--181 (2025). \doi{10.26464/epp2024058}

\bibitem{ref_shao_review}
Shao, M., et al.: Advances and challenges in solar flare prediction: A review. arXiv preprint arXiv:2511.20465 (2025). \doi{10.48550/arXiv.2511.20465}

\bibitem{ref_brown_walker}
Brown, D., Walker, A.: A semi-automatic method to measure the rotation of sunspots. Sol. Phys. \textbf{296}, 48 (2021). \doi{10.1007/s11207-021-01787-4}

\bibitem{ref_vit}
Dosovitskiy, A., Beyer, L., Kolesnikov, A., et al.: An image is worth 16x16 words: Transformers for image recognition at scale. In: International Conference on Learning Representations (2021)

\bibitem{ref_swin}
Liu, Z., Lin, Y., Cao, Y., et al.: Swin Transformer: Hierarchical vision transformer using shifted windows. In: Proceedings of the IEEE/CVF International Conference on Computer Vision (ICCV), pp. 10012--10022 (2021)

\bibitem{ref_convnext}
Liu, Z., Mao, H., Wu, C.-Y., Feichtenhofer, C., Darrell, T., Xie, S.: A ConvNet for the 2020s. In: Proceedings of the IEEE/CVF Conference on Computer Vision and Pattern Recognition (CVPR), pp. 11976--11986 (2022)

\bibitem{ref_mvit}
Fan, H., Xiong, B., Mangalam, K., Li, Y., Yan, Z., Malik, J., Feichtenhofer, C.: Multiscale Vision Transformers. In: Proceedings of the IEEE/CVF International Conference on Computer Vision (ICCV), pp. 6824--6835 (2021)

\bibitem{ref_mvitv2}
Li, Y., Wu, C.-Y., Fan, H., Mangalam, K., Xiong, B., Malik, J., Feichtenhofer, C.: MViTv2: Improved Multiscale Vision Transformers for Classification and Detection. In: Proceedings of the IEEE/CVF Conference on Computer Vision and Pattern Recognition (CVPR), pp. 4804--4814 (2022)


\bibitem{ref_long}
Pandey, C., Adeyeha, T., Hong, J., Angryk, R.A., Aydin, B.: Advancing solar flare prediction using deep learning with active region patches. In: Bifet, A., Krilavi\v{c}ius, T., Miliou, I., Nowaczyk, S. (eds.) Machine Learning and Knowledge Discovery in Databases. Applied Data Science Track. ECML PKDD 2024. Lecture Notes in Computer Science, vol. 14950. Springer, Cham (2024). \doi{10.1007/978-3-031-70381-2_4}


\bibitem{ref_soho}
SOHO project website, \url{https://soho.nascom.nasa.gov/}, last accessed 2026/03/29

\end{thebibliography}

\end{document}
